\subsection{\texorpdfstring{圆函数的导数；数 \(\pi\)}{圆函数的导数；数 π}}

我们在《一般拓扑》（Gen.~Top., VIII, p.~106）中定义过连续同态 \(x \mapsto \mathbf{e}(x)\) （加法群 \(\mathbf{R}\) 到绝对值为 1 的复数所成的乘法群 \(\mathbf{U}\) 上的）；这是一个以 1 为最小周期的周期函数，且 \(\mathbf{e}\left(\frac{1}{4}\right) = i\) 。我们知道（出处同上），\(\mathbf{R}\) 到 \(\mathbf{U}\) 上的每个连续同态都具有 \(x \mapsto \mathbf{e}(x / a)\) 的形式，并且令 \(\cos_a x = \mathcal{R}(\mathbf{e}(x / a))\) ，\(\sin_a x = \mathcal{I}(\mathbf{e}(x / a))\) （以 \(a\) 为底的三角函数或圆函数）；这些函数都是从 \(\mathbf{R}\) 到 \([-1, +1]\) 中的连续映射，以 \(a\) 为最小周期。我们有 \(\sin_a(x + a / 4) = \cos_a x\) ，\(\cos_a(x + a / 4) = -\sin_a x\) ，并且函数 \(\sin_a x\) 在区间 \([-a / 4, a / 4]\) 上递增。

命题 3. 函数 \(\mathbf{e}(x)\) 在 R 的每一点处有导数，等于 \(2\pi i\mathbf{e}(x)\) ，其中 \(\pi\) 是一个 \textgreater0 的常数。

现在，把 III, p.~51 的定理 1 应用于 E 是复数域 C 的情形，便得关系式 \(\mathbf{e}'(x) = \mathbf{e}'(0)\mathbf{e}(x)\) ；此外，由于 \(\mathbf{e}(x)\) 的欧氏范数为常数，\(\mathbf{e}'(x)\) 正交于 \(\mathbf{e}(x)\) （I, p.~7，例 3）；于是 \(\mathbf{e}'(0) = \alpha i\) ，其中 \(\alpha\) 为实数。由于 \(\sin_1 x\) 在 \([- \frac{1}{4}, \frac{1}{4}]\) 上递增，它在 \(x = 0\) 处的导数 \(\geqslant 0\) ，故 \(\alpha \geqslant 0\) ；又因 \(\mathbf{e}(x)\) 不是常数，故 \(\alpha > 0\) ；习惯上把这样得到的数 \(\alpha\) 记作 \(2\pi\) 。

在 §2（III, p.~23）中我们将说明如何计算数 \(\pi\) 的任意接近的近似值。

于是我们有公式

\[
\mathrm{D} \left(\mathbf {e} \left(\frac {x}{a}\right)\right) = \frac {2 \pi i}{a} \mathbf {e} \left(\frac {x}{a}\right).\tag{12}
\]

可以看出，当 \(a = 2\pi\) 时这个公式最简单；因此在分析学中人们只使用以 \(2\pi\) 为底的圆函数；我们约定在这些函数的记号中省略底；除非明确说明相反的情形，记号 \(\cos x\) ，\(\sin x\) 与 \(\tan x\) 分别表示 \(\cos_{2\pi}x\) ，\(\sin_{2\pi}x\) 与 \(\tan_{2\pi}x\) 。

在这些约定下，取 \(a = 2\pi\) ，公式 (12) 可写成

\[
\mathrm{D} (\cos x + i \sin x) = \cos \left(x + \frac {\pi}{2}\right) + i \sin \left(x + \frac {\pi}{2}\right),\tag{13}
\]

它等价于

\[
\mathrm{D} (\cos x) = - \sin x, \quad \mathrm{D} (\sin x) = \cos x,
\]

由此推出

\[
\mathrm{D} (\tan x) = 1 + \tan^ {2} x = \frac {1}{\cos^ {2} x}.\tag{15}
\]

除圆函数 \(\cos x\) ，\(\sin x\) 与 \(\tan x\) 之外，在数值计算中还使用三个辅助函数：余切、正割与余割，它们由下列公式定义

\[
\cot x = \frac {1}{\tan x}, \quad \sec x = \frac {1}{\cos x}, \quad \operatorname{cosec} x = \frac {1}{\sin x}.
\]

回想（Gen.~Top., VIII, p.~109），对应于底 \(2\pi\) 的角称为弧度。

